% Nội dung cắt khỏi mục 3.4 (chưa chạy). KHÔNG \input file này.
% - "Giới hạn diễn giải" -> Ch.5 (Trúc).
% - Bootstrap, SHAP, giá trị bắt được, chính sách theo giá trị đơn, ngân sách 10/20% -> Ch.6 hướng phát triển.

\section{Phương pháp đánh giá}\label{sec:eval}

\subsection{Chỉ số mô hình}

Do lớp dương chỉ chiếm khoảng một phần mười, độ chính xác (accuracy) không có ý nghĩa: một mô hình dự đoán mọi đơn đều tốt đã đạt khoảng 90\%. Chỉ số chính vì vậy là \textbf{PR-AUC}, diện tích dưới đường cong precision--recall, đo cả ở mức item và mức đơn (sau khi tổng hợp bằng \texttt{max-v1}). PR-AUC của mô hình ngẫu nhiên bằng đúng tỉ lệ dương, nên đây là mốc so sánh tự nhiên. ROC-AUC và precision, recall, F1 của lớp thiểu số tại các ngưỡng đã chọn được báo cáo bổ sung.

Độ tin cậy của xác suất được đo bằng Brier score:
\begin{equation}\label{eq:brier}
  \mathrm{BS} = \frac{1}{N}\sum_{j=1}^{N}\left(\hat{p}_j - y_j\right)^2
\end{equation}
trong đó $\hat p_j$ là xác suất dự đoán và $y_j \in \{0,1\}$ là nhãn thực tế. Brier score càng thấp càng tốt. Đường cong hiệu chỉnh (reliability diagram) được dùng để kiểm tra trực quan xác suất dự đoán có khớp với tỉ lệ thực tế ở từng khoảng điểm hay không.

\subsection{Hiệu quả ưu tiên}

Mục tiêu vận hành là giúp bộ phận chăm sóc khách hàng chọn đúng đơn khi chỉ có thể liên hệ một phần nhỏ. Với ngân sách $k \in \{5\%, 10\%, 20\%\}$ số đơn, nhóm đo:
\begin{itemize}
  \item \textbf{Recall@$k$}: tỉ lệ đơn nhận đánh giá thấp nằm trong top $k$ theo điểm rủi ro.
  \item \textbf{Giá trị bắt được@$k$}: tỉ lệ tổng giá trị của các đơn nhận đánh giá thấp nằm trong top $k$.
\end{itemize}
Danh sách theo mô hình được so sánh với hai chính sách đơn giản: chọn ngẫu nhiên (recall kỳ vọng bằng $k$) và chọn theo giá trị đơn từ cao xuống thấp. Chính sách thứ hai là đối chứng quan trọng vì nghiên cứu gần đây cho thấy xếp hạng theo giá trị đã biết có thể khó vượt qua bằng học máy khi năng lực xử lý hữu hạn~\cite{li2026valuesorting}. Đánh giá theo lợi ích ở mức ngân sách cụ thể, thay vì chỉ theo độ chính xác phân loại, cũng phù hợp với hướng tiếp cận dự đoán-tối ưu trong bài toán giữ chân khách hàng~\cite{gomezvargas2025}.

\subsection{Độ bất định và khả năng giải thích}

Khoảng tin cậy 95\% của các chỉ số được ước lượng bằng bootstrap 1\,000 lần với seed cố định. Việc lấy mẫu lại được thực hiện theo \emph{cụm đơn hàng}, không theo từng item, vì các item cùng đơn có chung nhãn và không độc lập với nhau.

Để giải thích mô hình, nhóm dùng giá trị SHAP ở mức toàn cục (nhóm đặc trưng nào ảnh hưởng nhiều nhất đến điểm rủi ro) và theo từng trường hợp (vì sao một đơn cụ thể bị xếp hạng cao). Giá trị SHAP mô tả cách mô hình sử dụng đặc trưng, không phải tác động nhân quả của đặc trưng lên sự hài lòng.

\subsection{Giới hạn diễn giải}

Kết quả của nghiên cứu cần được hiểu trong các giới hạn sau:
\begin{itemize}
  \item Không có đánh giá không có nghĩa là khách hài lòng, do quá trình tự chọn khi viết đánh giá~\cite{li2008}.
  \item Đánh giá thấp không đồng nghĩa với việc khách rời bỏ hoặc gây thiệt hại doanh thu.
  \item Điểm rủi ro cao không có nghĩa là khách sẽ phản hồi tốt với can thiệp. Thực nghiệm cho thấy nhóm khách rủi ro cao nhất và nhóm nhạy cảm nhất với can thiệp chỉ trùng nhau ở mức gần như ngẫu nhiên~\cite{ascarza2018}.
  \item Giá trị đơn bắt được trong danh sách ưu tiên là lợi ích xếp hạng, không phải doanh thu được cứu.
  \item Kết luận chỉ áp dụng cho giai đoạn dữ liệu Olist quan sát được, trừ khi được kiểm định trên dữ liệu bên ngoài.
\end{itemize}
